\documentclass[compress]{beamer}
\usepackage{ifthen,verbatim}

\newcommand{\isnote}{}
\xdefinecolor{lightyellow}{rgb}{1.,1.,0.25}
\xdefinecolor{darkblue}{rgb}{0.1,0.1,0.7}

%% Uncomment this to get annotations
%% \def\notes{\addtocounter{page}{-1}
%%            \renewcommand{\isnote}{*}
%% 	   \beamertemplateshadingbackground{lightyellow}{white}
%%            \begin{frame}
%%            \frametitle{Notes for the previous page (page \insertpagenumber)}
%%            \itemize}
%% \def\endnotes{\enditemize
%% 	      \end{frame}
%%               \beamertemplateshadingbackground{white}{white}
%%               \renewcommand{\isnote}{}}

%% Uncomment this to not get annotations
\def\notes{\comment}
\def\endnotes{\endcomment}

\setbeamertemplate{navigation symbols}{}
\setbeamertemplate{headline}{\mbox{ } \hfill
\begin{minipage}{5.5 cm}
\vspace{-0.75 cm} \small
\end{minipage} \hfill
\begin{minipage}{4.5 cm}
\vspace{-0.75 cm} \small
\begin{flushright}
\ifthenelse{\equal{\insertpagenumber}{1}}{}{Jim Pivarski \hspace{0.2 cm} \insertpagenumber\isnote/\pageref{numpages}}
\end{flushright}
\end{minipage}\mbox{\hspace{0.2 cm}}\includegraphics[height=1 cm]{../cmslogo} \hspace{0.1 cm} \includegraphics[height=1 cm]{../tamulogo} \hspace{0.01 cm} \vspace{-1.05 cm}}

\begin{document}
%% \begin{notes}
%% \item This is the annotated version of my talk.
%% \item If you want the version that I am presenting, download the one
%% labeled ``slides'' on Indico (or just ignore these yellow pages).
%% \item The annotated version is provided for extra detail and a written
%% record of comments that I intend to make orally.
%% \item Yellow notes refer to the content on the {\it previous} page.
%% \item All other slides are identical for the two versions.
%% \end{notes}

\small

\begin{frame}
\frametitle{Muon alignment planning}
\scriptsize
\vspace{0.5 cm}
\begin{tabular}{l l p{0.75\linewidth}}
Jim & done & \vspace{-1.5\baselineskip}\begin{itemize}
\item produced algorithm and software for R-T alignment
\item demonstrated ring-alignment (using monitoring plots only)
\item demonstrated beam-halo alignment (unmaintainable code)
\item documented almost everything in twiki format
\item got people interested in tracker global distortion issue
\end{itemize} \\
& to-do & \vspace{-1.5\baselineskip}\begin{itemize}
\item generalize and re-write beam-halo alignment procedure
\item document everything in the form of permanant notes
\end{itemize} \\

Aysen & done & \vspace{-1.5\baselineskip}\begin{itemize}
\item learned and improved R-T code
\end{itemize} \\
& to-do & \vspace{-1.5\baselineskip}\begin{itemize}
\item optimize R-T procedure for few pb$^{-1}$
\item start delivering constants: 2010 cosmic rays $\to$
  aligned barrel chambers and endcap ring-positions
\end{itemize} \\

Vadim & done & \vspace{-1.5\baselineskip}\begin{itemize}
\item findQualityFiles, Alignment Quality Monitor
\end{itemize} \\
& to-do & \vspace{-1.5\baselineskip}\begin{itemize}
\item hardware alignment case (?), more quality tests
\item formalize CSC ring-position procedure
\end{itemize} \\
\end{tabular}

%% \begin{itemize}  %%\setlength{\itemsep}{0.1 cm}

%% \end{itemize}
%% \hspace{-0.83 cm} \textcolor{darkblue}{\Large Outline2}
\end{frame}

\begin{frame}
\frametitle{2010 cosmic rays}

\begin{itemize}
\item Already collected: ``1 million cosmics'' (what that means
  depends on $p_T$ cut) ``with all subsystems participating''
  (essential).  Sometimes called ``CRAFT-10'' because it is similar in
  scope to CRAFT-09 and CRAFT-08.
\begin{itemize}
\item tracker alignment (prerequisite) already performed
\end{itemize}
\item First real alignment problem for Aysen: people need this for
  reconstruction of first TeV collisions data
\item Exact repeat of CRAFT-09 procedures: we've done it before, we can
  do it again, depending on statistics
\item Timescale: a few weeks
\item I'm planning a visit to A\&M next week, so that we can get
  started on this (should discuss optimal dates)
\end{itemize}
\end{frame}

\begin{frame}
\frametitle{Ring-position alignment}

\begin{columns}
\column{0.55\linewidth}
\includegraphics[height=\linewidth, angle=90]{ringfits_before_mem12.pdf}

\includegraphics[height=\linewidth, angle=90]{ringfits_after_mem12.pdf}

\column{0.4\linewidth}
\includegraphics[width=\linewidth]{tracker_disk_interpretation.pdf}
\end{columns}

\scriptsize
\begin{itemize}
\item Performed on monitoring plots (not in R-T framework), put positions in XML, produce SQLite file, and that becomes the alignment
\item Previously done by hand, but it can be naturally added to Vadim's diagnostics (easy to automate)
\end{itemize}
\end{frame}

%% \section*{First section}
%% \begin{frame}
%% \begin{center}
%% \Huge \textcolor{blue}{First section}
%% \end{center}
%% \end{frame}

\begin{frame}
\frametitle{What to do about CSC-Overlaps}

\includegraphics[width=0.9\linewidth]{occupancy_problems.pdf}

\begin{itemize}
\item CSC-Overlaps procedure depends on complete rings
\item Missing chamber/CFEB problems are not going to be resolved soon
\item Applying closure as a constraint only helps rings with one gap
  and is very sensitive to small problems (e.g.\ 8~mm error
  accumulated from 18$\times$500~$\mu$m effect)
\end{itemize}
\end{frame}

\begin{frame}
\frametitle{Generalizing procedure}

\includegraphics[width=0.5\linewidth]{incomplete_and_crosswise.pdf}

\begin{itemize}
\item Old procedure (left): nearest-neighbor constraints $m_{i,i+1}$ defines
  a static matrix whose inverse solves the alignment problem for all
  chambers simultaneously
\begin{itemize}
\item if not all connections are present, matrix would be non-invertable
\item filling in gaps by ``faking'' the missing residuals such that closure = 0
\end{itemize}

\item New procedure (right): constraints from any pair of chambers
  $m_{ij}$, matrix is built up dynamically from the constraints
\begin{itemize}
\item can treat constraints from residuals, photogrammetry, and maybe
  hardware alignment on an equal footing
\item requires a re-write of the module, which was needed anyway
\end{itemize}
\end{itemize}
\end{frame}           

\begin{frame}
\frametitle{Testing with 5-chamber example}

\scriptsize
\begin{minipage}{1.05\linewidth}

$\displaystyle \chi^2 = \sum_{\mbox{\scriptsize constraints}} \frac{(m_{ij} - A_i + A_j)^2}{{\sigma_{ij}}^2} + \lambda \left(\sum_i A_i\right)^2$

\vspace{0.3 cm}
\# format: ($i$, $j$, $m_{ij}$, $\sigma_{ij}$)

constraints = [(0, 1, 0.3, 1.), (1, 2, 0.1, 1.), (2, 3, -0.3, 1.), (3, 4, 0., 1.), (4, 0, -0.1, 1.),]

\vspace{0.3 cm}
result: $A_0 = 0.18$, $A_1 = -0.12$, $A_2 = -0.22$, $A_3 = 0.08$, $A_4 = 0.08$

\vspace{0.3 cm}
covariance: $\left(\begin{array}{r r r r r}
0.44 & 0.04 & -0.16 & -0.16 & 0.04 \\
0.04 & 0.44 & 0.04 & -0.16 & -0.16 \\
-0.16 & 0.04 & 0.44 & 0.04 & -0.16 \\
-0.16 & -0.16 & 0.04 & 0.44 & 0.04 \\
0.04 & -0.16 & -0.16 & 0.04 & 0.44 \\
\end{array}\right)$

\vspace{0.3 cm}
modes with uncorrelated uncertainties:

\vspace{0.1 cm}
\renewcommand{\arraystretch}{1.3}
\begin{tabular}{c | l}
$-0.45 A_{0} + -0.45 A_{1} + -0.45 A_{2} + -0.45 A_{3} + -0.45 A_{4}$ & $\propto \sqrt{1/\lambda}$ (meaningless) \\
$0.63 A_{0} + 0.25 A_{1} + -0.48 A_{2} + -0.54 A_{3} + 0.14 A_{4}$ & 0.60 (weaker modes) \\
$0.6 A_{1} + 0.37 A_{2} + -0.37 A_{3} + -0.6 A_{4}$ & 0.60 \\
$0.63 A_{0} + -0.38 A_{1} + -0.01 A_{2} + 0.4 A_{3} + -0.64 A_{4}$ & 0.37 (stronger modes) \\ 
$-0.39 A_{1} + 0.64 A_{2} + -0.64 A_{3} + 0.39 A_{4}$ & 0.37 \\
\end{tabular}

\vspace{0.3 cm}
Minuit gives the same results as linear algebra under many different configurations (but sometimes only in the first 3--4 digits for complicated cases)
\end{minipage}
\end{frame}

\begin{frame}
\frametitle{Planned documents}

\begin{itemize}
\item CSC Overlaps alignment procedure
\begin{itemize}
\item written at the same time as the new code
\item touches on every aspect of alignment theory in a simple example; will be a great introduction\ldots
\end{itemize}

\item Reference-Target alignment procedure and infrastructure
\begin{itemize}
\item including the XML tools, ring-alignment procedure
\end{itemize}
\end{itemize}

\vspace{0.2 cm}
\hspace{-0.83 cm} \textcolor{darkblue}{\Large Summary of procedures and software products}

\scriptsize
\vspace{0.1 cm}
\begin{itemize}
\item Reference-Target algorithm: chamber-by-chamber alignment for DTs and CSCs (AlignmentProducer module)
\item CSC-Overlaps algorithm: chambers-within-rings for CSCs (AlignmentProducer)
\item Scripts for submitting and collecting data from both of the above
\item CSC Ring alignment: ring positions relative to tracker (simple procedure performed on monitoring plots)
\item Monitoring plots in CommonAlignmentMonitor
\item XML tools for creating and analyzing muon system geometries
\item Alignment Quality Monitor browser for organizing study of monitoring plots (and eventually) differences between muon system geometries
\end{itemize}

\label{numpages}
\end{frame}

\end{document}
